\section{Why only the closed, grounded loop learns}
\label{sec:mechanism}

The categorical gap of Section~\ref{sec:results} follows from how actor--critic reinforcement learning assigns credit, and from what each stream can and cannot see. The actor improves its policy by following an advantage---roughly, ``is acting now better than the critic's estimate of the current state's value?''---and that signal is only useful if the critic is evaluating the same state the actor is acting from, and if the actor's input actually contains the evidence the action should depend on. The three configurations break, preserve, or misroute exactly these requirements.

\paragraph{V1: credit lands where it was earned.} In the closed, grounded loop the actor and critic operate over one shared deep memory $\mathbf{H}_2$, and the actor's own output writes that memory (Equation~\eqref{eq:arch-mem}). The critic therefore evaluates the same evolving state the actor occupies, and ``pressing now looks good'' is a statement about the state the actor actually inhabits. Credit lands on the representation that produced the action. Because the priority stream carries the raw image on its residual and reads both memories, the actor's input contains live visual evidence of the orientation change. The agent can therefore discover that declaring at change onset is rewarded, and can shape that policy by cue reliability and value---the covert-attention signature of Section~\ref{sec:results} emerges as the learned solution.

\paragraph{V2: credit is computed on the wrong representation.} When the streams are given private memories, the actor acts from $\mathbf{H}_{\mu}$ while the critic values $\mathbf{H}_{Q}$, and the two recurrent states drift apart over the episode. The advantage handed to the actor is then computed on a representation that is not the one the actor is acting from, so credit is systematically misattributed: an action that helped may be evaluated against a value estimate for a different internal state. Under such mis-specified feedback the policy falls into the only basin that is safe under \emph{any} value estimate---never act, never be penalised---and stays there. The always-wait collapse is the rational endpoint of an agent that cannot tell when acting helps. It is a learning failure, not a capacity failure: the network has the parameters and the pathways to express the task (V1 proves it), but with the loop open the teaching signal never reaches the right representation.

\paragraph{V3: the actor is fed a representation that cannot see the change.} Swapped routing preserves the coupling---value and policy are still bound through the shared $\mathbf{H}_2$, and credit assignment is clean---but the actor is now driven by the value stream, whose keys, values and residual are all the deep memory $\mathbf{H}_2$ rather than the image. That stream carries almost no live visual content: the image can only re-weight stored state, not inject new evidence. The decision is therefore made from a representation that cannot \emph{see} the orientation change as it happens. However well-posed the credit assignment, there is nothing change-related in the actor's input to assign credit to, and the policy again retreats to always-wait. V3 thus dissociates the two requirements: coupling alone is not enough; the stream that commands the action must be the grounded one.

\paragraph{Jointly necessary, individually insufficient.} The three outcomes form a clean logical decomposition. V2 has a grounded actor but an open loop and fails. V3 has a closed loop but a non-grounded actor and fails. Only V1 has both a closed loop and a grounded actor, and only V1 learns. Learning the task requires that the policy and value computations be coupled through a shared recurrent state \emph{and} that the image-grounded priority stream be the one that commands the policy---two conditions that a single structural edit can each remove, and whose removal is catastrophic rather than graded.
